


\exo 

On considère le point $A(2 ; -1)$ et le vecteur $\Coord{u}{2}{-3}$.

Le point $B(x_B ; y_B)$ est défini par $\vect{AB}=2\vect{u}$

\begin{enumerate}
\item Calculer les coordonnées de $2\vect{u}$
\item Exprimer les coordonnées de $\vect{AB}$ en fonction de $x_B$ et $y_B$
\item En déduire les coordonnées de $B$
\end{enumerate}

\medskip

\exo

On donne $A(-2 ; -3)$, $B(2 ; -1)$ et $C(1 ; 3)$

Déterminer les coordonnées du point $D$ tel que ABCD soit un parallélogramme.

\medskip

\exo

Soit $A(-2 ; 1)$, $B(2 ; 3)$, $C(3 ;-1)$ et $D(7 ;1)$

\begin{enumerate}
\item Calculer les coordonnées des milieux de $[AD]$ et de $[BC]$

Qu'en déduit-on ?
\item Proposer une autre méthode pour obtenir ce résultat.
\end{enumerate}
\exo
 
\begin{enumerate}
\item Placer les points $A(4 ; -2)$, $B(-1 ; 3.5)$ et $I(3;2)$
\item Construire les points $C$ et $D$ tels que ABCD soit un parallélogramme de centre $I$
\item Calculer les coordonnées de points $C$ et $D$
\end{enumerate}

\medskip

\exo 

Dans un repère orthonormal, on donne les points $A(-1 ; -1)$, $B(2 ; 3)$ et $C(4 ; -1)$.

\begin{enumerate}
\item Quelle est la nature du triangle ABC ?
\item Calculer le périmètre du triangle ABC.
\item Calculer l'aire du triangle ABC.
\end{enumerate}

\medskip

\exo 

On se place dans un repère orthonormé \Oij. L'unité de longueur est le centimètre.

On note $A$ le point de coordonnées $(-6 ; 8)$.

\begin{enumerate}
\item Déterminer les coordonnées du point $M$ milieu du segment $[OA]$.
\item Faire une figure et construire le cercle $\mathscr{C}$ de diamètre $[AO]$
\item Calculer la distance AO puis en déduire le rayon du cercle.
\item Le point $N(1;1)$ appartient-il au cercle $\mathscr{C}$ ?
\item Démontrer que le triangle NOA est rectangle en N. 
\end{enumerate} 

\medskip

\exo

On donne les points $A(-3;2)$, $B(0 ; -2)$ et $C(1 ; 5)$ dans un repère orthonormal.

\begin{enumerate}
\item Faire une figure que l'on complétera tout au long de l'exercice.
\item Calculer les longueurs AB et AC.
\item Déterminer les coordonnées du point D tel que ABDC soit un parallélogramme.
\item Quelle est la nature exacte de ABDC ? (Justifier)
\end{enumerate}

\medskip

\exo

\begin{enumerate}
\item Les vecteurs $\Coord{u}{9}{-6}$ et $\Coord{v}{-6}{4}$ sont-ils colinéaires ?
\item Les vecteurs $\Coord{u}{9}{-6}$ et $\Coord{v}{-6}{4}$ sont-ils colinéaires ?
\item Déterminer le nombre $x$ tel que $\Coord{u}{x}{3}$ et $\Coord{v}{2-x}{1}$ soient colinéaires ?
\item Les points $A(-5 ;1)$, $I(1 ;4)$ et $M(-1 ; 3)$ sont-ils alignés ?
\item On donne les points $A(-3 ; 3)$, $B(2;9)$, $C(-3 ;-4)$ et $D(6;4)$

Les droites $(AB)$ et $(CD)$ sont-elles parallèles ?
\end{enumerate}

\medskip

\exo

On considère $A(-2 ;2)$, $B(5;6)$, $C(4; -1)$ et $D(-3;-5)$.

\begin{enumerate}
\item Faire une figure.
\item Montrer que ABCD est un parallélogramme.
\item Déterminer les coordonnées du milieu $I$ du segment $[CD]$
\item Déterminer les coordonnées du point $M$ tel que $\vect{MC}=\dfrac{1}{3} \vect{AC}$.
\item Les vecteurs $\vect{BI}$ et $\vect{BM}$ sont-ils colinéaires ?

Que peut-on en déduire ?
\item On note $J$ le milieu du segment $[AB]$.

Montrer que les droites $(DJ)$ et $(BN)$ sont colinéaires.
\end{enumerate}

\medskip

\exo 

ABC est un triangle, I est le milieu de [AB], J est le point tel que $\overrightarrow{AJ}=\overrightarrow{AB}-2\overrightarrow{AC}$.
\begin{enumerate}
\item Faire une figure.
\item On choisit le repère $\left( A,\overrightarrow{AB},\overrightarrow{AC}\right) .$
\item Quelles sont les coordonnées des points I et J dans ce repère ?
\item Démontrez que les droites (AJ) et (CI) sont parallèles.
\end{enumerate}


